\section{Detect Cycle in a Directed Graph}
\label{sec:graph-cycle-directed}

\problemstatement{Given a directed graph, determine whether it
contains a cycle.}

\begin{patternbox}
Unlike the undirected case (the Graph BFS and DFS Problems chapter's Detect Cycle in an Undirected Graph section),
a single ``parent'' check is not enough here: in a \emph{directed}
graph, reaching an already-visited node is perfectly normal (two
different paths can legitimately converge on the same node without any
cycle). What signals a true cycle is reaching a node that is still
\textbf{on the current DFS path} -- tracked with a second array,
separate from \textit{visited}, marking only nodes currently ``in
progress.''
\end{patternbox}

\subsection*{Why \textit{visited} alone is insufficient here}

Consider $A \to B$, $A \to C$, $B \to D$, $C \to D$: node $D$ is
visited twice (once via $B$, once via $C$), but there is no cycle --
both are legitimate converging paths. The distinguishing feature of an
actual cycle is a \textbf{back edge}: an edge pointing to a node that
is an \textbf{ancestor} of the current node in the DFS recursion tree
(i.e., still actively on the call stack), not merely visited at some
earlier, already-concluded point in the traversal.

\subsection*{Approach 1: DFS with a Path-Tracking Array}

\begin{lstlisting}
#include <bits/stdc++.h>
using namespace std;

bool dfsHasCycle(int node, vector<vector<int>>& adj, vector<bool>& visited,
                  vector<bool>& pathVisited) {
    visited[node] = true;
    pathVisited[node] = true;

    for (int neighbor : adj[node]) {
        if (!visited[neighbor]) {
            if (dfsHasCycle(neighbor, adj, visited, pathVisited)) return true;
        } else if (pathVisited[neighbor]) {
            return true; // back edge to a node still on the current path
        }
    }
    pathVisited[node] = false; // backtrack: leaving this node's path
    return false;
}

bool hasCycleDFS(int n, vector<vector<int>>& adj) {
    vector<bool> visited(n, false), pathVisited(n, false);
    for (int i = 0; i < n; i++) {
        if (!visited[i] && dfsHasCycle(i, adj, visited, pathVisited)) {
            return true;
        }
    }
    return false;
}

int main() {
    int n = 3;
    vector<vector<int>> adj(n);
    adj[0] = {1}; adj[1] = {2}; adj[2] = {0};

    cout << (hasCycleDFS(n, adj) ? "true" : "false") << endl; // true
    return 0;
}
\end{lstlisting}

\begin{complexitybox}[Approach 1 Complexity]
\textbf{Time.} $O(V+E)$.
\textbf{Space.} $O(V)$ for \textit{visited}, \textit{pathVisited}, and
the recursion stack.
\end{complexitybox}

\subsection*{Approach 2: Kahn's Algorithm (BFS)}

As noted at the close of Section~\ref{sec:graph-topo-sort-kahn}, if
Kahn's algorithm processes fewer than $n$ nodes before its queue
empties, the unprocessed nodes are necessarily trapped in a cycle (a
cycle's nodes can never reach in-degree $0$, since each waits on
another node within the same cycle).

\begin{lstlisting}
#include <bits/stdc++.h>
using namespace std;

bool hasCycleKahn(int n, vector<vector<int>>& adj) {
    vector<int> inDegree(n, 0);
    for (int u = 0; u < n; u++) {
        for (int v : adj[u]) inDegree[v]++;
    }

    queue<int> q;
    for (int i = 0; i < n; i++) {
        if (inDegree[i] == 0) q.push(i);
    }

    int processed = 0;
    while (!q.empty()) {
        int node = q.front();
        q.pop();
        processed++;

        for (int neighbor : adj[node]) {
            inDegree[neighbor]--;
            if (inDegree[neighbor] == 0) q.push(neighbor);
        }
    }
    return processed != n; // fewer than n processed => cycle exists
}

int main() {
    int n = 3;
    vector<vector<int>> adj(n);
    adj[0] = {1}; adj[1] = {2}; adj[2] = {0};

    cout << (hasCycleKahn(n, adj) ? "true" : "false") << endl; // true
    return 0;
}
\end{lstlisting}

\subsection*{Dry run}

\dryrun{Take the $3$-cycle $0{\to}1{\to}2{\to}0$.}

\textbf{DFS}: $\textit{dfs}(0)$: visited$[0]=$true, pathVisited$[0]=$true.
Neighbor $1$: unvisited, recurse. $\textit{dfs}(1)$: visited,
pathVisited both true for $1$. Neighbor $2$: unvisited, recurse.
$\textit{dfs}(2)$: visited, pathVisited both true for $2$. Neighbor
$0$: \textbf{already visited, and pathVisited$[0]=$true} (still on the
current path) -- cycle detected, return \texttt{true}.

\textbf{Kahn's}: in-degrees all equal $1$ (each node has exactly one
incoming edge from the cycle). No node starts with in-degree $0$, so
the queue is empty from the start; $\textit{processed}=0 \neq 3=n$ --
cycle detected immediately, without even examining a single edge.

\begin{complexitybox}[Approach 2 Complexity]
\textbf{Time.} $O(V+E)$.
\textbf{Space.} $O(V)$ for in-degrees and the queue.
\end{complexitybox}

\begin{takeawaybox}
Directed-graph cycle detection needs to distinguish ``visited at some
point'' from ``visited and still active on this path'' -- the
\textit{pathVisited} array (cleared on backtrack) is what undirected
cycle detection's simpler parent check generalizes into once edges have
direction. Kahn's algorithm sidesteps this distinction entirely by
checking a global count instead, which is part of why it is often
preferred for this specific check.
\end{takeawaybox}
